\section{Discussion}
\label{sec:discussion}

\subsection{What the hypothesis got right and wrong}

We set out to test whether null results in misalignment training are diagnostic of robust, redundant or bottlenecked causal pathways, and whether a high \msp indicates hidden redundancy or entrenchment.

\para{Right: nulls are informative, and the machinery extracts that information.} Equivalence-tested bounds, a specification lattice, and a minimum-perturbation counterfactual all work as designed. Two induction cells here are licensed as genuine nulls at $\tau$ --- more than any null in the source literature establishes.

\para{Wrong, or at least unsupported: \msp as a property of the model.} The framing assumes \msp measures the model's causal structure. What this run shows is that the elicitation-side \msp is dominated by a property of the \emph{evaluation}: a single generic jailbreak flip. Before \msp can diagnose redundancy or entrenchment, it has to be made robust to that, which is what \mspa attempts and what a larger sample would let it deliver.

\para{Contradicted on the mechanism: redundancy.} Insofar as we could measure it, a high \msp is not accompanied by high representational redundancy. Effective ranks of $5.83$--$7.59$ against a random null of $14.9$ say these finetunes engage a concentrated, low-dimensional mechanism. That supports the bottleneck limb of the hypothesis over the redundancy limb --- though with the benign control missing, even this is provisional.

\subsection{The finding that generalises beyond this study}

The base-model result stands independently of every scope cut, because it is a comparison within a single fully-judged $16$-cell lattice at the largest sample sizes in the study. Its implication is concrete and immediately actionable: \textbf{papers that overturn a null by flipping an evaluation-side knob need to show the same flip does \emph{not} do it to an untrained model}. On our numbers a persona nudge alone takes an ordinary instruction-tuned model from $0.0\%$ to $93.8\%$ misalignment. Any ``the misalignment was there all along'' claim built on that kind of flip is not identified.

We want to be careful about the scope of this claim. It does not say prior elicitation results are wrong. \citet{conditional2026} use domain-cued triggers rather than a generic persona nudge, and our own lattice reproduces exactly that asymmetry: the \qset flips are the only ones that separate the released organism from the base model, while \persona separates nothing. The claim is that the \emph{class} of eval-side flips is heterogeneous, that some members of it are pure jailbreaks, and that only a base-model contrast tells you which is which.

\subsection{Limitations}
\label{sec:limitations}

\para{The judge is not \gptfouro, and this cuts a specific way.} Absolute rates are not comparable to published numbers, and $\tau$ is calibrated internally against a base-model floor and a released-organism ceiling; a different judge would move $\tau$ and could move \msp. The coherence judge also partly conflates content extremity with incoherence --- in our pre-run self-test a fluent but extreme response scored $24.9$ on coherence despite being grammatical --- which, because the outcome requires coherence $>50$, would delete the most egregious responses from the numerator and bias a study about nulls toward nulls. In the actual run this did not materialise (\secref{sec:results-checks}), and we report both gated and ungated rates so readers can verify rather than take it on trust. We did not run a second judge, so we cannot bound how much \msp moves under judge choice; this is the single most valuable follow-up.

\para{Scale, and what did not get run.} The compute budget was binding and was measured rather than assumed. Judge throughput did not batch away, and LoRA training collapsed from a benchmarked $1991$ tok/s to ${\approx}630$ tok/s per worker under two concurrent workers. The resulting cuts: seeds per lattice order were reduced to $3/2/1/1$, so higher-order cells rest on one training run each and \citet{overtrained2026} are explicit that single-seed results are unstable; organisms were trained on $1600$ examples against ${\approx}7000$ for the released adapters, making our organisms weaker by construction, which again biases the induction arm toward nulls; $8$ of $15$ organisms were trained; $43$ of $64$ cells were judged, with several cells at $n=3$--$4$ where nothing short of a near-total effect is detectable. \textbf{No benign-finetune control exists}, which is the most damaging single omission: it means the projection onto $d_{\mathrm{mis}}$ cannot be separated from generic finetuning movement, and the elicitation arm has no ``finetuned but harmless'' reference point. \mspi, the interaction decomposition and the mediation model are reported as not identified, not as nulls.

\para{Identification assumptions.} \msp is a counterfactual quantity, and following \citet{position2026identification} we state what it identifies under. (i) \emph{Lattice closure}: \msp is a minimum over the lattice we built. A factor we omitted --- optimizer choice, base-model family, dataset domain, sampling temperature --- could have crossed $\tau$ alone, which would make the true \msp smaller. \textbf{Every \msp value here is an upper bound}, and ``unreachable'' means unreachable \emph{within this lattice}, never impossible. (ii) \emph{Binary levels}: a dose between the two endpoints (say a $25\%$ mix) could cross $\tau$ where neither endpoint does, so the Hamming distance is only as meaningful as the level choices. (iii) \emph{Baseline dependence}: \msp is defined relative to $s^0$; a different null baseline is a different question. (iv) \emph{No interference between arms}: we treat training and evaluation flips as composable and measure the train $\times$ eval interaction at one pre-registered evaluation specification, not across the full $2^7$ product lattice, which was not affordable.

\para{Construct validity.} Misalignment as operationalised --- a judge's $0$--$100$ alignment score on eight free-form questions --- is a narrow proxy. A model can be dangerous in ways this instrument cannot see (deceptive alignment, situational awareness, long-horizon behaviour), and can score badly here for reasons that are not misalignment (an edgy persona, refusal-adjacent phrasing). The question sets are small, so we use cell counts rather than question-averaged rates throughout.

\para{Contamination and generality.} The neutral questions come from the public \texttt{first\_plot\_questions} set and may appear in post-training data for \qwenseven; the domain-cued set is newly written for this study, and notably it is the newly written set, not the public one, that carries the organism-attributable signal. Results are for one base model at one scale and one misalignment domain. \citet{overtrained2026} and \citet{evilspectra2026} both report large model-family effects, so cross-model generality of these \msp values is untested.

\para{What would invalidate these results.} A second judge producing a materially different $\tau$-ordering of the lattice cells; higher seed counts revealing that a cell we call crossing is a seed artifact, or that a licensed null is not; a single omitted factor crossing $\tau$ alone from the same baseline, collapsing a reported \msp of $2$ or $3$ to $1$; or the coherence gate being shown to remove genuinely misaligned-and-fluent responses often enough to flip a null classification.

\para{Relation to minimal-specification results elsewhere.} The closest analogue outside this literature is \citet{poisoning2025}, who find that ${\approx}250$ poisoned documents compromise models from 600M to 13B regardless of dataset size --- a minimum-specification result along a single count axis. \msp generalises that shape to a multi-factor lattice. In the other direction, \citet{sleeper2024} show backdoored behaviour persisting through SFT, RLHF and adversarial training, which is the large-\msp regime under safety-training flips that our induction arm was designed to detect and did not have the power to measure.

\subsection{Broader implications}

Three practices follow directly from these numbers. \textbf{Report an equivalence bound, not a point estimate}: a null should be published as ``the one-sided $95\%$ upper bound on the misalignment rate is $X\%$,'' together with the $n$ required to make that bound informative. On our numbers a null cell needs tens of responses to exclude even a $5\%$ rate, and many published nulls are reported at sample sizes that cannot exclude the effect they claim is absent. \textbf{Always run the base model through the same elicitation lattice}, because without it ``the misalignment was there all along'' and ``we jailbroke a normal model'' are indistinguishable. \textbf{Do not screen one factor at a time}, because where the effect lives in an interaction, single-flip screening reports a null for every factor individually while the joint flip crosses the threshold.

We also note the dual-use character of the persona result. It is a measurement of how easily an ordinary instruction-tuned open-weight model can be pushed into misaligned output by a prompt prefix, which is information a safety evaluation needs and which an attacker does not need from us --- the specific prefix class is already published \citep{persona2025nudge}. We report it because a safety literature that cannot tell jailbreaks from installed pathways will systematically over-credit its own mitigations.
